% ============================================================
% Appendix --- Problem sheet index
%
% This appendix exists because the restructure dissolved the official problem sheets: each
% problem now sits with the material it tests, so a sheet's problems can be spread over
% several chapters. The tables below put the sheet view back, and they are also where
% Corsin's colour-coded "Recommended exercises" pages are preserved --- those were
% per-sheet artefacts with nowhere else to live once the \exercisesheet{N} headings went.
%
% Built up during the migration, one sheet per stage.
% ============================================================
\chapter{Problem Sheet Index}
\label{ch:sheet_index}

% Transition: Claude Opus 5 (Medium)
The problem sheets are not reproduced as blocks anywhere in this document. Every problem is
quoted verbatim from the official sheet and placed next to the material it tests, with its
solution in the \emph{Solutions} section of the chapter that hosts it. That is the right
arrangement for reading the notes and the wrong one for answering \qt{what is on sheet 7?}, so
this appendix restores the sheet view.

All statements are quoted from \texttt{exercises/ExN\_Analysis2\_eng.pdf}; attribution:
Prof.\ Joaquim Serra, D-MATH, ETH Z\"urich.

\begin{notation}[Corsin's colour key]
\label{not:corsin_priority_key}
From sheet 4 onwards Corsin opens each file with a colour-coded \textit{Recommended exercises}
page. The key, as given there: blue \textbf{= important}, orange \textbf{= semi-important},
red \textbf{= optional}. The official sheets additionally mark harder problems with \textbf{(*)}.
Each problem carries its tag in its own title throughout this document; the tables below collect
them, together with Corsin's own comments.
\end{notation}

\begin{ainote}
Sheets 1--3 have no priority page --- Corsin's \textit{Recommended exercises} pages start with
sheet 4 --- so no problem on them is singled out, and all of them are reproduced and solved.
\end{ainote}

\section{Sheets 1--3}
\label{sec:sheet_index_1_3}

\begin{center}
\begin{tabular}{@{}l l l@{}}
\textbf{Problem} & \textbf{Topic} & \textbf{Now in} \\
\hline
1.1 & Recap questions          & \cref{sec:analysis1_recap} \\
1.2 & From equations to drawings & \cref{sec:visualizing_sets} \\
1.3 & From drawings to equations & \cref{sec:visualizing_sets} \\
1.4 & Young's inequality       & \cref{sec:youngs_inequality} \\[4pt]
2.1 & Metric spaces            & \cref{sec:metric_spaces} \\
2.2 & Multiple choice (closure) & \cref{sec:unit_balls} \\
2.3 & Multiple choice (distance from a set) & \cref{sec:unit_balls} \\
2.4 & Boundary, interior       & \cref{sec:unit_balls} \\
2.5 & Product of metric spaces & \cref{sec:metric_spaces} \\
2.6 & Continuity of the distance & \cref{sec:continuity} \\
2.7 & Continuity of the composition & \cref{sec:continuity} \\[4pt]
3.1 & Lipschitz or not         & \cref{sec:lipschitz_continuity} \\
3.2 & Problems at the origin   & \cref{sec:continuity} \\
3.3 & Open, closed, complete and compact & \cref{sec:why_compactness_matters} \\
3.4 & Multiple choice (open, complete, compact) & \cref{sec:why_compactness_matters} \\
3.5 & Multiple choice (fixed points) & \cref{sec:banach_fixed_point} \\
3.6 & Multiple choice (uniform continuity) & \cref{sec:lipschitz_continuity} \\
3.7 & Cauchy--Schwarz          & \cref{sec:cauchy_schwarz} \\
3.8 & Inner product and norm   & \cref{sec:cauchy_schwarz} \\
3.9 & All norms are equivalent & \cref{sec:cauchy_schwarz} \\
3.10 & Hilbert--Schmidt norm   & \cref{sec:cauchy_schwarz} \\
3.11 & (*) Uniform continuity and moduli & \cref{sec:lipschitz_continuity} \\
\end{tabular}
\end{center}

\section{Sheet 4}
\label{sec:sheet_index_4}

% Source: Corsin Nick/Class Notes/Week 4.pdf, p. 1
% Corsin's "Recommended exercises" page, preserved here: it was a per-sheet artefact and had
% nowhere else to go once the \exercisesheet{4} heading was dissolved.
\begin{center}
\begin{tabular}{@{}l p{0.22\linewidth} p{0.34\linewidth} l@{}}
\textbf{Problem} & \textbf{Priority} & \textbf{Corsin's note} & \textbf{Now in} \\
\hline
4.1 & important & ``Use the chain rule'' & \cref{sec:jacobi_matrix} \\[3pt]
4.2 & important & ``Recall that path connected implies connected'' & \cref{sec:path_connectedness} \\[3pt]
4.3 & semi-important --- parts 1, 2, 3, 5; \textbf{parts 4 and 6 important} &
  ``For 3., use Cauchy--Schwarz in $\mathbb{R}^n$ applied to convenient vectors.'' &
  \cref{sec:cauchy_schwarz} \\[3pt]
4.4 & optional & --- & \cref{sec:cauchy_schwarz} \\[3pt]
4.5 & optional & --- & \cref{sec:mean_value_theorem} \\[3pt]
4.6 & important --- parts 1, 2, 3 & --- & \cref{sec:chain_rule_along_path} \\
\end{tabular}
\end{center}

\begin{ainote}
Problems \textbf{4.4} and \textbf{4.5} appear in this document for the first time. They had been
transcribed into \texttt{content/exercise-sheets/week-04.tex}, but that file was never
\texttt{\textbackslash input} by anything, so neither problem had ever reached the built PDF and
neither carried a solution. Both are now placed and solved.
\end{ainote}
